\begin{enumerate}

\item Identificar si la sección está abierta, cerrada unicelular o bicelular; dibujar su geometría y espesores.

\item Para cada célula cerrada imponer equilibrio del par y compatibilidad del giro; el flujo común es la diferencia entre circulaciones.

\item Comprobar que el par exterior coincide con el que transmiten los flujos, si la sección es cerrada.

\item Calcular el giro total con theta = M\_x L/(G J), o con la compatibilidad de las celdas.

\item Las magnitudes de la solución se han contrastado mediante cálculo analítico y el motor estructural de Studium cuando éste ofrece esa comprobación.

\item En las ecuaciones bicelulares, $P_1=2w_1+h$ y $P_2=2w_2+h$ son los recorridos exteriores de cada célula; $t_e$ es el espesor exterior y $t_c$ el compartido; $r=q_2/q_1$.

\end{enumerate}

\noindent\textbf{Operaciones y resultados intermedios.}

\begin{enumerate}

\item Área de célula 1.

\[\fitmath{y_{1}=70\cdot 60\approx 4200\,\mathrm{mm^{2}}}\]

\item Área de célula 2.

\[\fitmath{y_{2}=70\cdot 60\approx 4200\,\mathrm{mm^{2}}}\]

\item Coeficiente a de compatibilidad.

\[\fitmath{y_{3}=\frac{\frac{200}{2}+\frac{60}{2}}{70\cdot 60}+\frac{60}{2\cdot 70\cdot 60}\approx 0.03809524\,\mathrm{1/mm^{2}}}\]

\item Coeficiente b de compatibilidad.

\[\fitmath{y_{4}=\frac{\frac{200}{2}+\frac{60}{2}}{70\cdot 60}+\frac{60}{2\cdot 70\cdot 60}\approx 0.03809524\,\mathrm{1/mm^{2}}}\]

\item Relación de flujos q2/q1 por compatibilidad.

\[\fitmath{y_{5}=\frac{\frac{\frac{200}{2}+\frac{60}{2}}{70\cdot 60}+\frac{60}{2\cdot 70\cdot 60}}{\frac{\frac{200}{2}+\frac{60}{2}}{70\cdot 60}+\frac{60}{2\cdot 70\cdot 60}}\approx 1\,}\]

\item Flujo circulante de la primera célula por equilibrio.

\[\fitmath{y_{6}=\frac{1680000}{2\cdot 70\cdot 60+2\cdot 70\cdot 60\cdot \frac{\frac{\frac{200}{2}+\frac{60}{2}}{70\cdot 60}+\frac{60}{2\cdot 70\cdot 60}}{\frac{\frac{200}{2}+\frac{60}{2}}{70\cdot 60}+\frac{60}{2\cdot 70\cdot 60}}}\approx 100\,\mathrm{N/mm}}\]

\item Flujo circulante de la segunda célula.

\[\fitmath{y_{7}=\frac{1680000}{2\cdot 70\cdot 60+2\cdot 70\cdot 60\cdot \frac{\frac{\frac{200}{2}+\frac{60}{2}}{70\cdot 60}+\frac{60}{2\cdot 70\cdot 60}}{\frac{\frac{200}{2}+\frac{60}{2}}{70\cdot 60}+\frac{60}{2\cdot 70\cdot 60}}}\cdot \frac{\frac{\frac{200}{2}+\frac{60}{2}}{70\cdot 60}+\frac{60}{2\cdot 70\cdot 60}}{\frac{\frac{200}{2}+\frac{60}{2}}{70\cdot 60}+\frac{60}{2\cdot 70\cdot 60}}\approx 100\,\mathrm{N/mm}}\]

\item Flujo neto de la pared compartida.

\[\fitmath{y_{8}=\frac{1680000}{2\cdot 70\cdot 60+2\cdot 70\cdot 60\cdot \frac{\frac{\frac{200}{2}+\frac{60}{2}}{70\cdot 60}+\frac{60}{2\cdot 70\cdot 60}}{\frac{\frac{200}{2}+\frac{60}{2}}{70\cdot 60}+\frac{60}{2\cdot 70\cdot 60}}}-\frac{1680000}{2\cdot 70\cdot 60+2\cdot 70\cdot 60\cdot \frac{\frac{\frac{200}{2}+\frac{60}{2}}{70\cdot 60}+\frac{60}{2\cdot 70\cdot 60}}{\frac{\frac{200}{2}+\frac{60}{2}}{70\cdot 60}+\frac{60}{2\cdot 70\cdot 60}}}\cdot \frac{\frac{\frac{200}{2}+\frac{60}{2}}{70\cdot 60}+\frac{60}{2\cdot 70\cdot 60}}{\frac{\frac{200}{2}+\frac{60}{2}}{70\cdot 60}+\frac{60}{2\cdot 70\cdot 60}}\approx 6.946121\times10^{-164}\,\mathrm{N/mm}}\]

\item Giro unitario compatible de ambas células.

\[\fitmath{y_{9}=\frac{\frac{\frac{1680000}{2\cdot 70\cdot 60+2\cdot 70\cdot 60\cdot \frac{\frac{\frac{200}{2}+\frac{60}{2}}{70\cdot 60}+\frac{60}{2\cdot 70\cdot 60}}{\frac{\frac{200}{2}+\frac{60}{2}}{70\cdot 60}+\frac{60}{2\cdot 70\cdot 60}}}\cdot 200}{2}+\frac{\left(\frac{1680000}{2\cdot 70\cdot 60+2\cdot 70\cdot 60\cdot \frac{\frac{\frac{200}{2}+\frac{60}{2}}{70\cdot 60}+\frac{60}{2\cdot 70\cdot 60}}{\frac{\frac{200}{2}+\frac{60}{2}}{70\cdot 60}+\frac{60}{2\cdot 70\cdot 60}}}-\frac{1680000}{2\cdot 70\cdot 60+2\cdot 70\cdot 60\cdot \frac{\frac{\frac{200}{2}+\frac{60}{2}}{70\cdot 60}+\frac{60}{2\cdot 70\cdot 60}}{\frac{\frac{200}{2}+\frac{60}{2}}{70\cdot 60}+\frac{60}{2\cdot 70\cdot 60}}}\cdot \frac{\frac{\frac{200}{2}+\frac{60}{2}}{70\cdot 60}+\frac{60}{2\cdot 70\cdot 60}}{\frac{\frac{200}{2}+\frac{60}{2}}{70\cdot 60}+\frac{60}{2\cdot 70\cdot 60}}\right)\cdot 60}{2}}{2\cdot 26955\cdot 70\cdot 60}\approx 4.416532\times10^{-5}\,\mathrm{rad/mm}}\]

\item Tensión tangencial máxima en las paredes.

\[\fitmath{y_{10}=\frac{\frac{1680000}{2\cdot 70\cdot 60+2\cdot 70\cdot 60\cdot \frac{\frac{\frac{200}{2}+\frac{60}{2}}{70\cdot 60}+\frac{60}{2\cdot 70\cdot 60}}{\frac{\frac{200}{2}+\frac{60}{2}}{70\cdot 60}+\frac{60}{2\cdot 70\cdot 60}}}}{2}\approx 50\,\mathrm{MPa}}\]

\end{enumerate}
